\section{Phase-aware memory layouts}
\label{sec:layout}

A $\Gamma$-invariant factor base has $2n$ times more points than
columns. Every table a pipeline keeps, the factor base itself, the
table of decomposed roots that a meet-in-the-middle probe looks up, and
the pair-sum table that a two-summand descent probes, can be stored
per orbit instead of per point, at the price of canonicalizing every
query. This section records the layouts the implementations use, the
bit budgets that make the phase free to store, and two observations
about canonical keys that are easy to get wrong.

\subsection{Representatives on disk}

The factor base is stored as one $x$-coordinate per signed orbit. For
the $n=83$ instance of Section~\ref{sec:bench-canon} the base of
$30\,977\,592$ usable points is $186\,612$ representatives, and
$30\,977\,592/186\,612=166=2n$ exactly, which is Lemma~\ref{lem:free}
checked on the data: every orbit is full. Each representative is $16$
bytes (the polynomial-basis $x$, little-endian), so the base is
$2.99$~MB instead of the $496$~MB that the points would occupy at the
same width.\art{experiments/koblitz-n83-index-design-20260929/adaptive_all_orientations_20261004/all83_ordinary_manifest.json}
Nothing else needs to be stored: the sign is recovered from the
canonical $x$ by the chooser, and the phase is recomputed, not stored,
because it is a function of the query, not of the base.

\subsection{The packed root record}
\label{sec:record}

The meet-in-the-middle decomposition at $n=83$ indexes the roots of
the summation polynomial $S_3$ over sampled pairs of base points and
then probes the index with the roots obtained from the target. The
index record is $18$ bytes:
\[
  \underbrace{\hat x}_{\text{11 bytes: exact 83-bit canonical key}}
  \;\|\;
  \underbrace{\text{offset}}_{\text{42 bits}}\,
  \underbrace{\text{root}}_{\text{1 bit}}\,
  \underbrace{k}_{\text{7 bits}}
  \quad(\text{7-byte descriptor}),
\]
where the offset addresses the satisfiable pair-state vector, the root
bit says which of the two $S_3$ roots this is, and $k<83$ is the
canonical shift of Definition~\ref{def:sf}.\art{experiments/koblitz-n83-index-design-20260929/packed_index_e1_canonical_20261009/native_packed_index.rs, encode}
The descriptor has $56$ bits and uses $50$; the phase costs nothing
beyond the bits that would otherwise pad the record. A probe
canonicalizes the partner root to $(\hat x',k')$, looks up $\hat x'$,
and recovers the relative Frobenius power as $(k+n-k')\bmod n$, which
is the exponent of $\tau$ that carries the stored root onto the
query. That is the only place the phase is consumed before relation
reconstruction, and it is an addition modulo $n$.

Records are sorted into $2^{22}$ buckets by a $22$-bit \emph{mixed}
prefix of the key (counting scatter, then a short sort inside each
bucket), with a directory of $32$-bit offsets. For $4\times10^6$
records the table is
$4\,000\,000\times18+2^{22}\times4+4=88\,777\,220$ bytes, against
$192.0$~MiB for the open-addressed hash table it replaced, a $2.27\times$
reduction in table bytes at the same record count (the first layout,
a $20$-bit raw prefix with $80\,388\,616$ bytes, had been
$2.50\times$).\art{experiments/koblitz-n83-index-design-20260929/packed_index_e1_prefix_20261007/README.md}
Since $4\times10^6<2^{22}$ an offset needs $22$ bits, and the ten
spare bits of each directory word hold an $8$-bit key fingerprint and
a singleton flag; a fingerprint mismatch on a singleton bucket rejects
the query before any record is loaded. On the frozen probe this cut
full-key comparisons from $9\,324\,648$ to $4\,943\,112$, with
$2\,912\,395$ queries rejected by the tag
alone.\art{experiments/koblitz-n83-index-design-20260929/packed_index_e1_tag_20261009/README.md}

\subsection{Canonical keys are skewed}
\label{sec:skew}

A canonical key is a \emph{minimum} over $n$ rotations, so its leading
bits are far from uniform. By Lemma~\ref{lem:zerogap} it begins with
$Z$ zeros where $Z$ is the longest cyclic zero run of the word, and for
a uniform $n$-bit word $\Pr[Z\ge z]\approx1-\exp(-n/2^{z+1})$, so at
$n=83$ the top five bits are zero with probability about $0.73$ and
the top six with probability about $0.48$. A directory indexed by the
raw leading bits of the key therefore collapses into a small fraction
of its buckets. This was observed twice, independently, before it was
understood: a folded pair table that bucketed on the key's high bits
``piled into the first few buckets'', invisible on misses and
expensive on every hit, and the first packed index at $n=83$ performed
$58\,251\,024$ full-key comparisons with a raw $20$-bit prefix against
$25\,265\,242$ with a hashed one.\art{crypto:docs/ic/runs/koblitz-orbit-fold-20260913.json, a_bug_the_tag_exposed; experiments/koblitz-n83-index-design-20260929/packed_index_e1_hashed_prefix_20261008/README.md}
The rule that follows is simple and we state it as one: \emph{never
index a canonical key by its leading bits; hash it, or index by a
mixed prefix.} The same skew is a feature for compression. A canonical
key has $\lfloor\log_2 n\rfloor$ or more leading zeros with constant
probability, so a variable-length encoding of the key saves a few
bits per record; the implementations do not exploit this, and we
record it as an open engineering item.

\subsection{The folded pair table and the orbit tag}
\label{sec:pairtable}

The two-summand descent asks, for a residual $R$, whether
$R=P_i+P_j$ for base points $P_i,P_j$, by looking $x(R)$ up in a table
of all pair sums. Since the base is $\Gamma$-invariant the set of pair
sums is $\Gamma$-invariant too ($\tau(P+Q)=\tau P+\tau Q$ and
$-(P+Q)=(-P)+(-Q)$), so the table needs one key per orbit of sums.
Three tiers exist in the implementation: \emph{full} (every pair, with
both summands stored), \emph{compact} (every pair, summands
recovered by a scan) and \emph{folded} (one key per orbit of sums,
canonical key as above). The folded tier stores $2n$ times fewer
entries and therefore affords a base $\sqrt{2n}$ times wider in the
same memory, which divides the descent's expected $2r/|F|^2$ probes
by $2n$; it pays with a canonicalization on every probe.

Two details make the folded tier usable rather than merely small.

\paragraph{Enumerating each orbit of sums once.}
Walking $i$ over one representative of each signed orbit and $j$ over
the whole base reaches every orbit of ordered pairs (carry $P_i$ to
its representative by some $g\in\Gamma$; then $(gP_i,gP_j)$ is
enumerated), so no sum orbit is missed. But a sum orbit is reached
twice, once from each summand's orbit, and the first implementation
stored it twice, which made the fold worth $n$ rather than $2n$ and
made an ``equal memory'' comparison not equal. The row filter
``row $\alpha$ keeps $j$ only when $\orb(j)\ge\alpha$'' stores each sum
orbit once and keeps the covering
property.\art{crypto:docs/ic/runs/koblitz-orbit-fold-20260913.json, the_row_filter}

\paragraph{Recovery in $O(n)$, not $O(|F|)$.}
A folded entry says only that the orbit of $R$ contains a sum
$P_{\mathrm{rep}(\alpha)}+P_j$. Recovering the actual summands of $R$,
which the relation needs, would mean scanning the base. Storing the
orbit index $\alpha$ of one summand in the high half of the entry's
``rest'' word (the bits the bucket hash does not already determine)
turns the scan into a walk over the $2n$ points of orbit $\alpha$:
$R=g^{-1}P_{\mathrm{rep}(\alpha)}+g^{-1}P_j$ with the first summand in
orbit $\alpha$, so one tries the $2n$ members of that orbit and
subtracts. Measured at $n=61$ the recovery fell from $8.66$~ms to
$0.037$~ms per decomposition, from the majority of a decomposition to
$1.4\%$ of it, at no cost in table
bytes.\art{crypto:docs/ic/runs/koblitz-orbit-fold-20260913.json, the_orbit_tag}
This is the second place the phase is consumed: the orbit tag is the
``which column'' half of $\can(P)$, and the walk over the orbit is the
``which phase'' half, computed on demand.

\paragraph{The degenerate key.}
$O$ is the sum of every $\pm$ pair, so the orbit of $O$, the single
element $O$, is a key under which every orbit representative stores an
entry. A lookup there returns every orbit at once. It is bounded
($14.02$~ms for the $88\,816$ pairs it legitimately returns at the
measured width) once the two membership tests inside it are sorts
rather than scans, and it is a reminder that Remark~\ref{rem:short}
applies to sums as well as to points.
